\section{Interior critical points}\label{sec:interior}

Section \ref{sec:compactness} converts a failure of $\Gtz(m,k)$ into an
extremal failure: one at a design minimizing the objective, rather than
at an arbitrary design.  This section extracts what such a design
satisfies.  The constraint set is a smooth manifold, and its tangent
space is computed first.  The objective is a maximum of least-eigenvalue
functions, and its first-order behaviour is that of a nonsmooth maximum.
The two combine into a finite system of algebraic identities, and
everything drawn from that system afterwards is algebra.

A design at $(m,k)$ over $\R$ is written as a pair $(g,t)$ with
$g=(g_1,\dots,g_m)\in(\R^k)^m$ and $t=(t_1,\dots,t_m)$.  Put
\begin{equation}\label{eq:designobjective}
  F(g,t)\;:=\;\max_{\lvert C\rvert=k}\lmin\bigl(\SC\bigr),
  \qquad \SC=\sum_{c\in C}g_c^{\vphantom{\mathsf{T}}}g_c\tp.
\end{equation}
By Definition~\ref{def:design} and the eigenvalue reading recorded in
Lemma~\ref{lem:tiemax}, a $k$-subset dominates exactly when
$\lmin(\SC)\ge1$, so $\Gtz(m,k)$ asserts $F\ge1$ on every design, and a
failure of $\Gtz(m,k)$ is a design with $F<1$.  The value of $F$ is never
zero.

\begin{lemma}\label{lem:fpos}
$F(g,t)>0$ for every design.
\end{lemma}

\begin{proof}
By \eqref{eq:parseval} the vectors span $\R^k$, so some $k$-subset $C$ is
a basis; equivalently, $\sum_{\lvert C\rvert=k}\pi_C=1$ in
\eqref{eq:dpp} while $\pi_C>0$ only for bases.  For a basis and $x\ne0$,
$x\tp\SC x=\sum_{c\in C}\langle g_c,x\rangle^2>0$, so $\lmin(\SC)>0$.
\end{proof}

The first-order theory below is applied at a design where $F$ is
minimal, and one statement is needed to place such a design where the
variational principle put its own.

\begin{definition}\label{def:extremalcex}
The cell $(m,k)$ \emph{has an extremal counterexample} if $F$ attains a
local minimum, over the designs of size $m$ and rank $k$, at a design
whose value is below $1$.
\end{definition}

Suppose $\Gtz(m,k)$ fails.  The designs with $F<1$ then form a nonempty
set, open in the design manifold, and a local minimum of $F$ on that set
is a local minimum over all designs, since $F\ge1$ off it.  By
Lemma~\ref{lem:chart} those designs are exactly the designs whose chart
points carry no dominating $k$-subset, and, granting $\Gtz(m-1,k)$ and
$\Gtz(m-1,k-1)$, Corollary~\ref{cor:interiorcex} places the whole of that
locus in the compactification of \S\ref{sec:compactness} strictly inside
the open weight simplex.  What is not automatic is that the infimum of $F$ over the
locus is attained rather than approached at the boundary of the
compactification, where by Theorem~\ref{thm:boundary} some subset
dominates.  Theorem~\ref{thm:variational} settles precisely that question
for the chart objective $G$; and by Remark~\ref{rem:objectives} $G$ and
$F$ agree in sign but are related by a congruence which is not a scalar,
so a minimizer of one need not minimize the other.
Definition~\ref{def:extremalcex} names what the passage requires, and it
is carried as a hypothesis in Theorem~\ref{thm:rankthree} and nowhere
else.

\subsection{The design manifold}\label{sec:manifold}

Fix $m$ and $k$, write $\mathbb{S}^k$ for the space of real symmetric
$k\times k$ matrices, and let
\begin{equation}\label{eq:constraints}
  \Phi(g,t)\;:=\;\sum_{c}t_c\,g_c^{\vphantom{\mathsf{T}}}g_c\tp-I_k
  \;\in\;\mathbb{S}^{k},
  \qquad
  \sigma(t)\;:=\;\sum_c t_c-1\;\in\;\R .
\end{equation}
The designs are the points of $\{\Phi=0\}\cap\{\sigma=0\}$ satisfying the
open condition $t\pd0$.  Both maps are polynomial, with differentials
\begin{equation}\label{eq:conGrad}
  \Phi'(g,t)[\dot g,\dot t]
  =\sum_c\Bigl(\dot t_c\,g_c^{\vphantom{\mathsf{T}}}g_c\tp
      +t_c\bigl(\dot g_c^{\vphantom{\mathsf{T}}}g_c\tp
                +g_c^{\vphantom{\mathsf{T}}}\dot g_c\tp\bigr)\Bigr),
  \qquad
  \sigma'(t)[\dot t]=\sum_c\dot t_c .
\end{equation}

\begin{lemma}\label{lem:tangent}
At every design the pair $(\Phi',\sigma')$ is surjective onto
$\mathbb{S}^k\times\R$.  Hence the designs form a smooth manifold of
dimension $mk+m-\tfrac12k(k+1)-1$, whose tangent space $\tang$ at $(g,t)$
consists of the pairs $(\dot g,\dot t)$ annihilated by
\eqref{eq:conGrad}.
\end{lemma}

\begin{proof}
Let $(Y,a)\in\mathbb{S}^k\times\R$ be given.  Choose $\dot t$ with
$\sum_c\dot t_c=a$, put
$Y':=Y-\sum_c\dot t_c\,g_c^{\vphantom{\mathsf{T}}}g_c\tp$, and take
$\dot g_c:=\tfrac12Y'g_c$.  Then by \eqref{eq:parseval}
\[
  \sum_c t_c\bigl(\dot g_c^{\vphantom{\mathsf{T}}}g_c\tp
                 +g_c^{\vphantom{\mathsf{T}}}\dot g_c\tp\bigr)
  =\tfrac12Y'\sum_c t_c\,g_c^{\vphantom{\mathsf{T}}}g_c\tp
   +\tfrac12\Bigl(\sum_c t_c\,g_c^{\vphantom{\mathsf{T}}}g_c\tp\Bigr)Y'
  =Y' ,
\]
so $\Phi'[\dot g,\dot t]=Y$ and $\sigma'[\dot t]=a$.  The dimension count
and the description of $\tang$ are the submersion theorem.
\end{proof}

For an antisymmetric $X$ the infinitesimal rotation
\begin{equation}\label{eq:orbit}
  \dot g_c\;=\;Xg_c\quad(1\le c\le m),\qquad \dot t=0
\end{equation}
lies in $\tang$, since
$\sum_c t_c\bigl(Xg_c^{\vphantom{\mathsf{T}}}g_c\tp
+g_c^{\vphantom{\mathsf{T}}}g_c\tp X\tp\bigr)=X-X=0$ by \eqref{eq:parseval}; and
$X\mapsto(Xg_c)_c$ is injective, because the $g_c$ span $\R^k$ and a
matrix killing a spanning set is zero.  Write $\mathcal{O}\subseteq \tang$
for the resulting subspace, of dimension $\tfrac12k(k-1)$.  The
objective \eqref{eq:designobjective} is constant along $\mathcal{O}$.

\subsection{The first-order system}\label{sec:kkt}

The objective is a maximum of finitely many least-eigenvalue functions of
matrices depending polynomially on $g$.  It is locally Lipschitz, not
differentiable, and not Clarke regular: each $\lmin(\SC)$ is a
\emph{concave} function of $\SC$, and for a maximum of functions that are
not regular the subdifferential is contained in --- and in general not
equal to --- the convex hull of the subdifferentials of the active pieces
\cite[Prop.~2.3.12]{Clarke1990}.  The inclusion runs in the direction a
necessary condition needs.

The function $\lmin$ is concave on $\mathbb{S}^k$, and its
superdifferential at $A$ is the \emph{spectraplex}
\begin{equation}\label{eq:spectraplex}
\begin{aligned}
  \partial\lmin(A)\;&:=\;\bigl\{\,\Theta\psd0:\ \tr\Theta=1,\
    \operatorname{ran}\Theta\subseteq\ker\bigl(A-\lmin(A)I\bigr)\,\bigr\}\\
  &\;=\;\operatorname{conv}\bigl\{uu\tp:\ Au=\lmin(A)u,\
    \lVert u\rVert=1\bigr\}
\end{aligned}
\end{equation}
\cite[\S2]{Overton1988},
\cite[Vol.~I, Ch.~VI]{HiriartUrrutyLemarechal1993}.  The Clarke subdifferential of a locally Lipschitz function composed
with a smooth map is, in turn, contained in the image of the subdifferential under the adjoint of
the derivative \cite[Thm.~2.3.10]{Clarke1990}.  Applied to
$\lmin\circ\SC$, and using that $\Theta$ is symmetric so that
$\langle\Theta,\dot g_c^{\vphantom{\mathsf{T}}}g_c\tp
+g_c^{\vphantom{\mathsf{T}}}\dot g_c\tp\rangle
=2\langle\Theta g_c,\dot g_c\rangle$, these give
\begin{equation}\label{eq:pieceSub}
  \partial\bigl(\lmin\circ \SC\bigr)(g,t)
  \;\subseteq\;
  \Bigl\{\,\phi_{C,\Theta}:(\dot g,\dot t)\mapsto
    2\sum_{c\in C}\langle\Theta g_c,\dot g_c\rangle
    \ :\ \Theta\in\partial\lmin(\SC)\,\Bigr\}.
\end{equation}
The functionals $\phi_{C,\Theta}$ do not involve $\dot t$.

\begin{theorem}\label{thm:firstorder}
Let $(g,t)$ be a design at which $F$ attains a local minimum over the
designs.  Put $\cv:=F(g,t)$ and
$\mathcal{A}:=\{C:\lvert C\rvert=k,\ \lmin(\SC)=\cv\}$.  Then there are a
finite index set $A$, subsets $C_i\in\mathcal{A}$, reals $\lambda_i\ge0$
with $\sum_{i\in A}\lambda_i=1$, unit vectors $u_i$ with
$S_{C_i}u_i=\cv u_i$, a symmetric matrix $\Lambda$ and a scalar $\mu$ such
that
\begin{align}
  \sum_{i\,:\,c\in C_i}\lambda_i\langle u_i,g_c\rangle\,u_i
    &\;=\;t_c\,\Lambda g_c &&(1\le c\le m), \label{eq:statg}\\
  g_c\tp\Lambda\,g_c &\;=\;\mu &&(1\le c\le m). \label{eq:statt}
\end{align}
\end{theorem}

\begin{proof}
Compose $F$ with a smooth local parametrization of the manifold of
Lemma~\ref{lem:tangent}.  The composition has an unconstrained local
minimum, so zero lies in its Clarke subdifferential
\cite[Prop.~2.3.2]{Clarke1990}; by the chain rule
\cite[Thm.~2.3.10]{Clarke1990} there is $\psi\in\partial F(g,t)$
vanishing on $\tang$.  Inactive subsets remain inactive near $(g,t)$, so the
maximum rule \cite[Prop.~2.3.12]{Clarke1990} together with
\eqref{eq:pieceSub} exhibits $\psi$ as a convex combination
\begin{equation}\label{eq:zeta}
  \psi\;=\;\sum_{C\in\mathcal{A}}\lambda_C\,\phi_{C,\Theta_C},
  \qquad \lambda_C\ge0,\quad\sum_C\lambda_C=1,\quad
  \Theta_C\in\partial\lmin(\SC).
\end{equation}
Decompose each $\Theta_C$ spectrally,
$\Theta_C=\sum_jw_{C,j}u_{C,j}^{\vphantom{\mathsf{T}}}u_{C,j}\tp$ with
$w_{C,j}\ge0$, $\sum_jw_{C,j}=1$ and each $u_{C,j}$ a unit eigenvector of
$\SC$ at $\cv$; this is possible by \eqref{eq:spectraplex}.  Index by the
pairs $i=(C,j)$, setting $C_i:=C$, $u_i:=u_{C,j}$ and
$\lambda_i:=\lambda_Cw_{C,j}$.  The new multipliers are nonnegative and
sum to one, and \eqref{eq:zeta} becomes
\begin{equation}\label{eq:zetarankone}
  \psi(\dot g,\dot t)
  \;=\;2\sum_c\Bigl\langle
        \sum_{i\,:\,c\in C_i}\lambda_i\langle u_i,g_c\rangle\,u_i,\
        \dot g_c\Bigr\rangle .
\end{equation}
By Lemma~\ref{lem:tangent} the map $(\Phi',\sigma')$ is surjective with
kernel $\tang$, and $\psi$ vanishes on $\tang$; hence $\psi$ factors through
$(\Phi',\sigma')$.  That is, there are $\Lambda\in\mathbb{S}^k$ and
$\mu\in\R$ with
\begin{equation}\label{eq:lagrange}
  \psi(\dot g,\dot t)
  \;=\;\bigl\langle\Lambda,\ \Phi'(g,t)[\dot g,\dot t]\bigr\rangle
      \;-\;\mu\,\sigma'(t)[\dot t] ,
\end{equation}
the sign of $\mu$ being a normalization.  Expanding the right side by
\eqref{eq:conGrad}, and using
$\langle\Lambda,\dot g_c^{\vphantom{\mathsf{T}}}g_c\tp
+g_c^{\vphantom{\mathsf{T}}}\dot g_c\tp\rangle=2\langle\Lambda g_c,\dot g_c\rangle$
as before,
\[
  \bigl\langle\Lambda,\Phi'[\dot g,\dot t]\bigr\rangle
   -\mu\sum_c\dot t_c
  \;=\;
  2\sum_c t_c\langle\Lambda g_c,\dot g_c\rangle
  \;+\;\sum_c\dot t_c\bigl(g_c\tp\Lambda g_c-\mu\bigr).
\]
The coordinates $\dot g$ and $\dot t$ are independent, so comparison with
\eqref{eq:zetarankone} may be made term by term.  The coefficient of
$\dot g_c$ gives \eqref{eq:statg} after division by $2$, and the
coefficient of $\dot t_c$ gives \eqref{eq:statt}.
\end{proof}

The maximum rule in \eqref{eq:zeta} is only an inclusion, because the pieces
are concave and therefore not Clarke regular; the sharper first-order
statement available at such a point is the directional one, namely
$\max_{C\in\mathcal{A}}\lmin\bigl(Q_C\tp\dot S_CQ_C\bigr)\ge0$ for every
$(\dot g,\dot t)\in \tang$, where $Q_C$ is an orthonormal basis of the
$\lmin$-eigenspace of $\SC$ and $\dot S_C$ is the derivative of $\SC$
\cite{HiriartUrrutyYe1995,Danskin1967}.  Nothing below uses more than the multiplier form.  The infinitesimal
rotations \eqref{eq:orbit}, for their part, impose no condition: for $\Theta\in\partial\lmin(\SC)$ one has
$\Theta\SC=\cv\Theta=\SC\Theta$, so
\begin{equation}\label{eq:rotationkill}
  \phi_{C,\Theta}\bigl[(Xg_c)_c,0\bigr]
  =2\sum_{c\in C}\langle\Theta g_c,Xg_c\rangle
  =2\tr\bigl(\Theta X\SC\bigr)=2\cv\,\tr(\Theta X)=0,
\end{equation}
because $\Theta$ is symmetric and $X$ antisymmetric.

\subsection{Stationarity data}\label{sec:data}

The conclusion of Theorem~\ref{thm:firstorder} is a finite algebraic
object, and it is convenient to detach it from its provenance.

\begin{definition}\label{def:datum}
A \emph{stationarity datum} for a design $(g,t)$ at $(m,k)$, with value
$\cv\in\R$, consists of a finite index set $A$, $k$-subsets $C_i$
($i\in A$), reals $\lambda_i\ge0$ with $\sum_i\lambda_i=1$, vectors $u_i$
with $\lVert u_i\rVert=1$ and $S_{C_i}u_i=\cv u_i$, and a symmetric matrix
$\Lambda$, satisfying \eqref{eq:statg} and \eqref{eq:statt}.  Write
$A^{+}:=\{C_i:\lambda_i>0\}$ for the \emph{active family}, a set of
$k$-subsets.  The datum is \emph{admissible} if in addition
\begin{equation}\label{eq:argmax}
  \cv\;\ge\;\lmin(\SC)\qquad\text{for every }k\text{-subset }C ,
\end{equation}
which says $\cv\ge F(g,t)$.  It says $\cv=F(g,t)$ exactly when in addition
$\cv=\lmin(S_{C_i})$ for some $i$, which the definition does not require;
Theorem~\ref{thm:firstorder} produces data for which it holds.
\end{definition}

The subsets $C_i$ need not be distinct: a tight eigenvalue of
multiplicity $r$ contributes $r$ indices carrying one subset, which is
what the spectral decomposition in the proof of
Theorem~\ref{thm:firstorder} produces.  Nothing below assumes that
$i\mapsto C_i$ is injective, and this is why the active family is defined
as a set of subsets rather than as a set of indices.

Definition~\ref{def:datum} asks only that $\cv$ be \emph{an} eigenvalue of
each $S_{C_i}$.  It does not ask that $\cv=\lmin(S_{C_i})$, and it does not
ask for \eqref{eq:argmax}.  Theorem~\ref{thm:firstorder} supplies both;
the algebra of the next two subsections supplies neither, and
\S\ref{sec:reach} shows that the difference is not a formality.

\subsection{The quadric law and the coverage law}\label{sec:laws}

Pairing \eqref{eq:statg} with $g_c$ gives, for every $c$,
\begin{equation}\label{eq:contract}
  \sum_{i\,:\,c\in C_i}\lambda_i\langle u_i,g_c\rangle^2
  \;=\;t_c\,\bigl(g_c\tp\Lambda g_c\bigr) .
\end{equation}
Both laws of the next proposition come from this contraction.  Throughout we use the
polarized form of the subset sum,
\begin{equation}\label{eq:polarized}
  y\tp\SC x\;=\;\sum_{c\in C}\langle g_c,y\rangle\langle g_c,x\rangle ,
\end{equation}
immediate from \eqref{eq:dominate}.

\begin{proposition}[Lean]\label{prop:quadric}
Let a stationarity datum be given.  Then $\mu=\cv$, so that every vector
lies on one central quadric of the multiplier,
\begin{equation}\label{eq:quadric}
  g_c\tp\Lambda\,g_c\;=\;\cv\qquad(1\le c\le m);
\end{equation}
for every $c$ the coverage law
\begin{equation}\label{eq:coverage}
  \sum_{i\,:\,c\in C_i}\lambda_i\langle u_i,g_c\rangle^2\;=\;t_c\,\cv
\end{equation}
holds; the multiplier is determined by the remaining data,
\begin{equation}\label{eq:multiplier}
  \Lambda\;=\;\cv\sum_{i\in A}\lambda_i\,u_i^{\vphantom{\mathsf{T}}}u_i\tp;
\end{equation}
and consequently $\cv\ge0$, $\Lambda\psd0$ and $\tr\Lambda=\cv$.
\end{proposition}

\begin{proof}
Sum \eqref{eq:contract} over $c$.  On the right, \eqref{eq:statt} makes
$g_c\tp\Lambda g_c$ independent of $c$, so the right side is
$\mu\sum_ct_c=\mu$.  On the left, exchange the two summations; for fixed
$i$ the inner sum runs over the $k$ vectors of $C_i$ and equals, by
\eqref{eq:polarized} and $S_{C_i}u_i=\cv u_i$,
\[
  \sum_{c\in C_i}\langle u_i,g_c\rangle^2
  \;=\;u_i\tp S_{C_i}u_i\;=\;\cv\lVert u_i\rVert^2\;=\;\cv .
\]
Hence the left side is $\cv\sum_i\lambda_i=\cv$, so $\mu=\cv$.  Substituting
$\mu=\cv$ into \eqref{eq:statt} gives \eqref{eq:quadric}, and into
\eqref{eq:contract} gives \eqref{eq:coverage}.

For \eqref{eq:multiplier}, fix $x\in\R^k$ and resolve it through
\eqref{eq:parseval}: $x=\sum_ct_c\langle g_c,x\rangle g_c$.  Then
\begin{align*}
  \Lambda x
  &=\sum_c\langle g_c,x\rangle\,\bigl(t_c\Lambda g_c\bigr)
   =\sum_c\langle g_c,x\rangle
      \sum_{i\,:\,c\in C_i}\lambda_i\langle u_i,g_c\rangle\,u_i
  &&\text{by \eqref{eq:statg}}\\
  &=\sum_{i\in A}\lambda_i
      \Bigl(\sum_{c\in C_i}\langle g_c,u_i\rangle
        \langle g_c,x\rangle\Bigr)u_i
   =\sum_{i\in A}\lambda_i\bigl(u_i\tp S_{C_i}x\bigr)u_i
  &&\text{by \eqref{eq:polarized}}\\
  &=\cv\sum_{i\in A}\lambda_i\langle u_i,x\rangle\,u_i ,
\end{align*}
the last step by symmetry of $S_{C_i}$ together with
$S_{C_i}u_i=\cv u_i$.  A matrix is determined by its action, which is
\eqref{eq:multiplier}.  Finally
$\cv=u_i\tp S_{C_i}u_i=\sum_{c\in C_i}\langle g_c,u_i\rangle^2\ge0$ for any
$i$, and \eqref{eq:multiplier} then exhibits $\Lambda$ as a nonnegative
multiple of a convex combination of rank-one projectors, so
$\Lambda\psd0$ and $\tr\Lambda=\cv\sum_i\lambda_i=\cv$.
\end{proof}

The coverage law \eqref{eq:coverage} is the contraction
\eqref{eq:contract} rewritten by the quadric law, and the quadric law is
that same contraction summed over $c$.  The two are therefore not
independent constraints, and no count of the equations of the system may
treat them as such: the content of a stationarity datum is exactly
\eqref{eq:statg} and \eqref{eq:statt}.

Reading the quadric law through \eqref{eq:multiplier} removes the
restriction $c\in C_i$ from the sum.

\begin{lemma}[Lean]\label{lem:overlap}
Let a stationarity datum with $\cv\ne0$ be given.  Then
\begin{equation}\label{eq:overlapone}
  \sum_{i\in A}\lambda_i\langle u_i,g_c\rangle^2\;=\;1
  \qquad(1\le c\le m),
\end{equation}
and consequently
\begin{enumerate}
\item[\rm(i)] $\ell_c\ge1$ for every $c$;
\item[\rm(ii)] $t_c\cv\le1$ for every $c$, and hence $\cv\le m$;
\item[\rm(iii)] for every $k$-subset $C$, active or not, there is an
  index $i\in A$ with $u_i\tp\SC u_i\ge k$;
\item[\rm(iv)] for each $c$ there is $i$ with $\lambda_i>0$, $c\in C_i$
  and $\langle u_i,g_c\rangle\ne0$; in particular $A^{+}$ covers
  $\{1,\dots,m\}$.
\end{enumerate}
\end{lemma}

\begin{proof}
By \eqref{eq:multiplier},
$g_c\tp\Lambda g_c=\cv\sum_i\lambda_i\langle u_i,g_c\rangle^2$; comparison
with \eqref{eq:quadric} and cancellation of $\cv\ne0$ give
\eqref{eq:overlapone}.

(i) Cauchy--Schwarz and $\lVert u_i\rVert=1$ give
$\langle u_i,g_c\rangle^2\le\ell_c$; substituting in
\eqref{eq:overlapone} and using $\sum_i\lambda_i=1$ gives $1\le\ell_c$.

(ii) The sum in \eqref{eq:coverage} is a subsum of the sum in
\eqref{eq:overlapone} with nonnegative terms, so $t_c\cv\le1$.  Summing
over $c$ against $\sum_ct_c=1$ gives $\cv\le m$.

(iii) Sum \eqref{eq:overlapone} over the $k$ vectors of $C$ and use
\eqref{eq:polarized}: $\sum_i\lambda_i\,u_i\tp\SC u_i=k$.  A convex
combination does not exceed its largest term.

(iv) By \eqref{eq:coverage} the sum at $c$ equals $t_c\cv\ne0$, so one of
its terms is nonzero, and that forces $\lambda_i>0$, $c\in C_i$ and
$\langle u_i,g_c\rangle\ne0$ for the corresponding $i$.
\end{proof}

Part (iii) names a unit vector; it implies $\lmax(\SC)\ge k$ for every $k$-subset, but the vector is the
useful half.  Part (i) says every vector of a stationary design is heavy
in the weak sense; the strict inequality $\ell_c>1$ would require
excluding equality in Cauchy--Schwarz at every positively weighted index
simultaneously, and is not available.

\subsection{The tetrahedron}\label{sec:tetracalib}

Before drawing exclusions we verify the entire system on the design of
Example~\ref{ex:tetra}, where every constant in it can be displayed.

\begin{example}\label{ex:tetrakkt}
Take the tetrahedron: $g_c=\sqrt3\,u_c$ with $u_c$ the four unit vertex
vectors, $t_c=\tfrac14$, so that $\ell_c=3$ and
$\langle g_i,g_j\rangle=-1$ for $i\ne j$.  There are four triples, each
the complement of one index; for $C=\{0,1,2,3\}\setminus\{c_0\}$ one has
$\SC=4I_3-g_{c_0}^{\vphantom{\mathsf{T}}}g_{c_0}\tp$, with spectrum $\{1,4,4\}$.
Take
\[
  A=\{0,1,2,3\},\qquad C_i=\{0,1,2,3\}\setminus\{i\},\qquad
  \lambda_i=\tfrac14,\qquad
  u_i=\tfrac{1}{\sqrt3}\,g_i,\qquad \cv=1 .
\]
Each $u_i$ is a unit vector, and it is the tight eigenvector:
$\langle g_i,u_i\rangle=3/\sqrt3=\sqrt3$, so
$S_{C_i}u_i=4u_i-\sqrt3\,g_i=4u_i-3u_i=u_i$, the eigenvalue being
$\cv=1$.  The multiplier predicted by \eqref{eq:multiplier} is
\[
  \Lambda\;=\;1\cdot\sum_{i}\tfrac14\cdot\tfrac13\,
    g_i^{\vphantom{\mathsf{T}}}g_i\tp
  \;=\;\tfrac1{12}\sum_ig_i^{\vphantom{\mathsf{T}}}g_i\tp
  \;=\;\tfrac1{12}\cdot4I_3\;=\;\tfrac13I_3 ,
\]
using $\sum_ig_i^{\vphantom{\mathsf{T}}}g_i\tp=4I_3$, which is
\eqref{eq:parseval} at $t_i=\tfrac14$.  Then $\tr\Lambda=1=\cv$, as
Proposition~\ref{prop:quadric} requires, and the quadric law
\eqref{eq:quadric} reads $g_c\tp\Lambda g_c=\ell_c/3=1=\cv$.

For the coverage law, the triples containing $c$ are the three with
$i\ne c$, and for those
$\langle u_i,g_c\rangle=\langle g_i,g_c\rangle/\sqrt3=-1/\sqrt3$, so
\[
  \sum_{i\,:\,c\in C_i}\lambda_i\langle u_i,g_c\rangle^2
  \;=\;3\cdot\tfrac14\cdot\tfrac13\;=\;\tfrac14\;=\;t_c\cv .
\]
Identity \eqref{eq:overlapone} restores the missing index $i=c$, for
which $\langle u_c,g_c\rangle^2=3$, and indeed
$\tfrac14+\tfrac14\cdot3=1$.  For \eqref{eq:statg} itself, using
$\sum_ig_i=0$,
\[
  \sum_{i\,:\,c\in C_i}\lambda_i\langle u_i,g_c\rangle u_i
  \;=\;\sum_{i\ne c}\tfrac14\Bigl(-\tfrac1{\sqrt3}\Bigr)
       \tfrac{g_i}{\sqrt3}
  \;=\;-\tfrac1{12}\sum_{i\ne c}g_i
  \;=\;\tfrac{g_c}{12}
  \;=\;\tfrac14\cdot\tfrac13\,g_c
  \;=\;t_c\Lambda g_c .
\]
The datum is admissible: every triple has $\lmin=1$, so $F=1=\cv$.  The
bounds of Lemma~\ref{lem:overlap} read $\ell_c=3\ge1$,
$t_c\cv=\tfrac14\le1$ and $\cv=1\le4=m$.  Part (iii) of that lemma is not
realized at the tight index of the subset itself, since
$u_i\tp S_{C_i}u_i=\cv=1<3$; it is realized at any other index, for which
\[
  u_j\tp S_{C_i}u_j=4-\langle g_i,u_j\rangle^2
   =4-\tfrac{\langle g_i,g_j\rangle^2}{3}=4-\tfrac13=\tfrac{11}3\;\ge\;3
   \qquad(j\ne i).
\]
\end{example}

\subsection{Exclusions}\label{sec:exclusions}

Each exclusion below is a statement about the datum alone;
admissibility is nowhere used.

\begin{proposition}\label{prop:e1}
Let a stationarity datum with $\cv\ne0$ have $\Lambda=a\,ww\tp$ for
some unit vector $w$ and some $a\in\R$.  Then $\cv=k$.  Since
$\Lambda\psd0$, this covers every datum with $\rk\Lambda\le1$.
\end{proposition}

\begin{proof}
Taking traces in \eqref{eq:multiplier} gives $a=\tr\Lambda=\cv$.
Evaluating \eqref{eq:multiplier} as a quadratic form at an arbitrary $x$
gives $\cv\langle w,x\rangle^2=\cv\sum_i\lambda_i\langle u_i,x\rangle^2$, so
after cancelling $\cv\ne0$,
\begin{equation}\label{eq:rankonecollapse}
  \sum_{i\in A}\lambda_i\langle u_i,x\rangle^2
  \;=\;\langle w,x\rangle^2\qquad\text{for every }x .
\end{equation}
Put $x=w$.  Then $\sum_i\lambda_i\langle u_i,w\rangle^2=1
=\sum_i\lambda_i$, so
$\sum_i\lambda_i\bigl(1-\langle u_i,w\rangle^2\bigr)=0$; every summand is
nonnegative, because $\lVert u_i\rVert=\lVert w\rVert=1$ forces
$\langle u_i,w\rangle^2\le1$.  Hence every summand vanishes.  Fix $i$
with $\lambda_i>0$.  Then $\langle u_i,w\rangle^2=1$ and
\[
  \bigl\lVert u_i-\langle u_i,w\rangle w\bigr\rVert^2
  \;=\;1-2\langle u_i,w\rangle^2+\langle u_i,w\rangle^2\;=\;0 ,
\]
so $u_i=\pm w$.  Now read \eqref{eq:quadric} through $\Lambda=\cv\,ww\tp$:
it says $\cv\langle w,g_c\rangle^2=\cv$, so $\langle w,g_c\rangle^2=1$ for
every $c$.  Therefore, by \eqref{eq:polarized},
\[
  \cv\;=\;u_i\tp S_{C_i}u_i\;=\;w\tp S_{C_i}w
    \;=\;\sum_{c\in C_i}\langle w,g_c\rangle^2
    \;=\;\lvert C_i\rvert\;=\;k. \qedhere
\]
\end{proof}

\begin{corollary}\label{cor:e1}
No stationarity datum with $0<\cv<1$ has a multiplier of rank at most one.
\end{corollary}

\begin{proof}
Otherwise $\cv=k$ by Proposition~\ref{prop:e1}, and $\cv\ne0$ forces
$k\ge1$, whence $\cv\ge1$.
\end{proof}

The second exclusion concerns the tight directions.  Say that a datum has
\emph{independent tight support} if the only relation
$\sum_{i\in A}\lambda_i\kappa_i\,u_i=0$ with $\kappa:A\to\R$ has
$\lambda_i\kappa_i=0$ for every $i$; linear independence of the family
$(u_i)_{\lambda_i>0}$ implies this, and the stated form makes the
zero-multiplier indices cost nothing.

\begin{proposition}\label{prop:e2}
Let a stationarity datum with $\cv\ne0$ have independent tight support.
Then $t_c\cv=1$ for every $c$; hence $\cv=m$ and $t_c=1/m$ for every $c$.
In particular no stationarity datum with $0<\cv<1$ has independent tight
support: below the threshold the tight directions carried by positive
multipliers are always dependent.
\end{proposition}

\begin{proof}
Substituting \eqref{eq:multiplier} into the right side of
\eqref{eq:statg} turns it into a combination of the $u_i$,
\[
  t_c\Lambda g_c
  \;=\;t_c\cv\sum_{i\in A}\lambda_i\langle u_i,g_c\rangle\,u_i ,
\]
so subtraction gives the relation
\begin{equation}\label{eq:vanishing}
  \sum_{i\in A}\lambda_i\,\langle u_i,g_c\rangle\,
    \bigl(\,[\,c\in C_i\,]-t_c\cv\,\bigr)\,u_i\;=\;0 ,
\end{equation}
where $[\,\cdot\,]$ is $1$ when the condition holds and $0$ otherwise.
Independence kills every coefficient:
\[
  \lambda_i\,\langle u_i,g_c\rangle\,
  \bigl(\,[\,c\in C_i\,]-t_c\cv\,\bigr)\;=\;0\qquad(i\in A).
\]
Fix $c$.  By Lemma~\ref{lem:overlap}(iv) there is $i$ with
$\lambda_i>0$, $c\in C_i$ and $\langle u_i,g_c\rangle\ne0$; for that $i$
the last factor is $1-t_c\cv$, so $t_c\cv=1$.  Summing over $c$ against
$\sum_ct_c=1$ gives $\cv=m$, and then $t_c=1/m$.  A design has at least one
vector, so $\cv=m\ge1$.
\end{proof}

The conclusion identifies the datum rather than contradicting it: one
might expect independence to force $\langle u_i,g_c\rangle=0$, against
coverage, but the coefficient in \eqref{eq:vanishing} vanishes through
its last factor instead.

Counting gives a third exclusion: by Lemma~\ref{lem:overlap}(iv) the
active family covers $\{1,\dots,m\}$, and each of its members has $k$
elements.

\begin{proposition}\label{prop:count}
For a stationarity datum with $\cv\ne0$,
\begin{equation}\label{eq:count}
  m\;\le\;k\,\lvert A^{+}\rvert .
\end{equation}
If $k<m$ then $\lvert A^{+}\rvert\ge2$.  If $\lvert A^{+}\rvert=1$ then
$m=k$ and $\cv=k$.  If $m=2k$ and $\lvert A^{+}\rvert=2$, the two members
of $A^{+}$ are disjoint and partition $\{1,\dots,m\}$.
\end{proposition}

\begin{proof}
The union of the members of $A^{+}$ is all of $\{1,\dots,m\}$ and has at
most $k\lvert A^{+}\rvert$ elements, which is \eqref{eq:count}; the
second assertion is \eqref{eq:count} read backwards.  If $A^{+}=\{C\}$
then $m\le k$, while $m\ge k$ always, so $m=k$; moreover every index with
$\lambda_i>0$ carries the same subset $C$, so summing
\eqref{eq:overlapone} over $c\in C$ gives, as in
Lemma~\ref{lem:overlap}(iii), $\sum_i\lambda_iu_i\tp\SC u_i=k$, while
each such $u_i$ is a unit eigenvector of $\SC$ at $\cv$, whence $\cv=k$.  If
$m=2k$ and $A^{+}=\{C,C'\}$ then $C\cup C'$ has $2k$ elements and
$\lvert C\rvert+\lvert C'\rvert=2k$, so $C\cap C'=\emptyset$.
\end{proof}

\begin{corollary}[Lean]\label{cor:singleactive}
A stationarity datum with $\cv\ne0$ whose index set $A$ is a singleton has
$\cv=k$.  In particular no such datum has $0<\cv<1$.
\end{corollary}

\begin{proof}
A singleton $A=\{i_0\}$ has $\lambda_{i_0}=1>0$, so $A^{+}=\{C_{i_0}\}$
and Proposition~\ref{prop:count} gives $\cv=k$; and $\cv\ne0$ forces $k\ge1$,
whence $\cv\ge1$.
\end{proof}

\begin{corollary}\label{cor:e3}
Let a stationarity datum at $(6,3)$ have $0<\cv<1$.  Then either
$\lvert A^{+}\rvert\ge3$, or $A^{+}$ consists of two disjoint triples
partitioning $\{1,\dots,6\}$ and the tight directions carried by positive
multipliers are linearly dependent.  At $(7,3)$ one has
$\lvert A^{+}\rvert\ge3$ unconditionally.
\end{corollary}

\begin{proof}
At $(6,3)$, Proposition~\ref{prop:count} gives
$\lvert A^{+}\rvert\ge2$; if it is exactly $2$ then $m=2k$ forces the
partition, and Proposition~\ref{prop:e2} forbids independent tight
support at $\cv<1$.  At $(7,3)$, \eqref{eq:count} reads
$7\le3\lvert A^{+}\rvert$.
\end{proof}

Whether the partition branch at $(6,3)$ is inhabited is open.

A fourth exclusion follows from the same subsum comparison, and it acts
on the opposite end of the range: on families so large that a single
vector meets all of them.

\begin{proposition}[Lean]\label{prop:saturated}
Let a stationarity datum with $\cv\ne0$ be given, and suppose some index
$c$ lies in $C_i$ for every $i\in A$.  Then $t_c\cv=1$, and hence
$\cv\ge1$.
\end{proposition}

\begin{proof}
The sum in \eqref{eq:coverage} runs over $\{i:c\in C_i\}$ and the sum in
\eqref{eq:overlapone} over all of $A$.  Under the hypothesis the two
index sets coincide, so the inequality $t_c\cv\le1$ of
Lemma~\ref{lem:overlap}(ii) is an equality: $t_c\cv=1$.  Since
$t_c\le1$ this gives $\cv=1/t_c\ge1$.
\end{proof}

\begin{corollary}\label{cor:saturated}
At a datum with $0<\cv<1$, no index lies in every member of $A$; in
particular, for each index some active subset omits it.
\end{corollary}

Propositions~\ref{prop:e1}, \ref{prop:e2}, \ref{prop:count} and
\ref{prop:saturated} are the exclusions available at $0<\cv<1$.  Their
combined reach is measured in \S\ref{sec:census}.

\subsection{The reach of the system}\label{sec:reach}

The exclusions bite on small active families.  At the frontier size the
system is satisfiable at the threshold itself, with a large active
family.

\begin{example}\label{ex:splitseven}
Split the tetrahedron of Example~\ref{ex:tetra} at three of its four
vertices.  Let $h_0,\dots,h_3$ be the sign vectors $(1,1,1)$,
$(1,-1,-1)$, $(-1,1,-1)$, $(-1,-1,1)$, so that
$\sum_dh_d^{\vphantom{\mathsf{T}}}h_d\tp=4I_3$ and $\langle h_d,h_e\rangle=-1$
for $d\ne e$; these are the vectors of Example~\ref{ex:tetra} in
coordinates.  Give the seven vectors $g_0,\dots,g_6$ the directions
$h_0,h_1,h_2,h_3,h_0,h_1,h_2$ and the weights
$\tfrac18,\tfrac18,\tfrac18,\tfrac14,\tfrac18,\tfrac18,\tfrac18$.  Each
direction carries total weight $\tfrac14$, so \eqref{eq:parseval} holds.
This is one of the extremal designs of
Theorem~\ref{thm:tieseverywhere} at $(7,3)$, and $F=1$.

Call a triple \emph{rainbow} if its three vectors point in three distinct
directions.  There are $2^3=8$ rainbow triples on the directions
$h_0,h_1,h_2$, one for each choice of vector in each direction; and a
rainbow triple containing the unsplit vector $g_3$ uses two of the three
split directions, which gives $3\cdot2\cdot2=12$ more, twenty in all.
For a rainbow triple $C$ missing the direction $h_d$ one has
$\SC=4I_3-h_d^{\vphantom{\mathsf{T}}}h_d\tp$, exactly as in
Example~\ref{ex:tetrakkt}, so $u_C:=h_d/\sqrt3$ is a unit eigenvector at
$\cv=1$.  Give multiplier $\tfrac1{32}$ to the eight triples omitting $h_3$
and $\tfrac1{16}$ to the twelve containing $g_3$; these are positive and
$8\cdot\tfrac1{32}+12\cdot\tfrac1{16}=\tfrac14+\tfrac34=1$.  Each of the
four directions is the missed one for total multiplier mass $\tfrac14$,
so \eqref{eq:multiplier} gives
\[
  \Lambda\;=\;1\cdot\sum_{d}\tfrac14\cdot\tfrac13\,
    h_d^{\vphantom{\mathsf{T}}}h_d\tp\;=\;\tfrac1{12}\cdot4I_3
    \;=\;\tfrac13I_3 ,
\]
whence $\tr\Lambda=1=\cv$ and $g_c\tp\Lambda g_c=3/3=1=\cv$, which is
\eqref{eq:quadric}.  A triple containing $g_c$ misses a direction other
than that of $g_c$, so $\langle u_C,g_c\rangle^2=\tfrac13$ there, and the
coverage law reduces to $\tfrac13\sum_{C\ni c}\lambda_C=t_c$.  For $c=3$
the twelve triples of multiplier $\tfrac1{16}$ all contain $g_3$, giving
$\tfrac13\cdot\tfrac{12}{16}=\tfrac14=t_3$.  For $c=0$ the triples
containing $g_0$ are four on the directions $h_0,h_1,h_2$, of multiplier
$\tfrac1{32}$, and two each on $h_0,h_1,h_3$ and on $h_0,h_2,h_3$, of
multiplier $\tfrac1{16}$, giving
\[
  \tfrac13\Bigl(\tfrac4{32}+\tfrac2{16}+\tfrac2{16}\Bigr)
  =\tfrac13\cdot\tfrac38=\tfrac18=t_0 .
\]
The datum is admissible, and $\lvert A^{+}\rvert=20$.
\end{example}

Example~\ref{ex:splitseven} puts twenty triples in the active family at the cell to which rank
three reduces.  It does not conflict
with Proposition~\ref{prop:carath} below, which asserts only that the
multipliers \emph{may be taken} supported on at most nineteen indices.

The next statement is the reason the conjectures of \S\ref{sec:census}
carry an admissibility hypothesis.

\begin{theorem}\label{thm:reach}
There is a design at $(4,2)$ carrying a stationarity datum with
$\cv=2-\sqrt2<1$.  That design satisfies $\Gtz(4,2)$ with margin: some
$2$-subset $C$ has $\SC-I=I\pd0$, so $F=2$.  Consequently the
implication ``a stationarity datum with $\cv\ne0$ has $\cv\ge1$'' is false,
and no strengthening of the identities \eqref{eq:statg} and
\eqref{eq:statt} can make it true.
\end{theorem}

\begin{proof}
Take
\[
  g_0=(\sqrt2,0),\quad g_1=(-1,1),\quad g_2=(1,1),\quad g_3=(0,\sqrt2),
  \qquad t_c=\tfrac14 .
\]
Then $\ell_c=2$ for every $c$ and
\[
  \sum_ct_c\,g_c^{\vphantom{\mathsf{T}}}g_c\tp
  =\tfrac14\left[
    \begin{pmatrix}2&0\\0&0\end{pmatrix}
   +\begin{pmatrix}1&-1\\-1&1\end{pmatrix}
   +\begin{pmatrix}1&1\\1&1\end{pmatrix}
   +\begin{pmatrix}0&0\\0&2\end{pmatrix}\right]
  =\tfrac14\begin{pmatrix}4&0\\0&4\end{pmatrix}=I_2 ,
\]
so this is a design.  Take the four pairs $C_0=\{0,1\}$, $C_1=\{0,2\}$,
$C_2=\{1,3\}$, $C_3=\{2,3\}$, each with multiplier $\tfrac14$, and
$\Lambda:=\tfrac{2-\sqrt2}{2}I_2$.  Now
\[
  S_{C_0}=\begin{pmatrix}3&-1\\-1&1\end{pmatrix},
  \qquad \tr S_{C_0}=4,\qquad\det S_{C_0}=2,
\]
so the eigenvalues of $S_{C_0}$ are $2\pm\sqrt2$, and the same holds for
$C_1$, $C_2$, $C_3$ by symmetry.  Put $\cv:=2-\sqrt2$.  Unit eigenvectors
at $\cv$ are the normalizations $u_i:=r_i/\lVert r_i\rVert$ of
\[
  r_0=(1,1+\sqrt2),\quad r_1=(1,-(1+\sqrt2)),\quad
  r_2=(1,\sqrt2-1),\quad r_3=(1,1-\sqrt2),
\]
with $\lVert r_0\rVert^2=\lVert r_1\rVert^2=4+2\sqrt2$ and
$\lVert r_2\rVert^2=\lVert r_3\rVert^2=4-2\sqrt2$; for instance
\[
  \bigl(S_{C_0}-\cv I\bigr)r_0
  =\begin{pmatrix}1+\sqrt2&-1\\-1&\sqrt2-1\end{pmatrix}
   \begin{pmatrix}1\\1+\sqrt2\end{pmatrix}
  =\begin{pmatrix}(1+\sqrt2)-(1+\sqrt2)\\
                  -1+(\sqrt2-1)(1+\sqrt2)\end{pmatrix}=0 .
\]
Condition \eqref{eq:statt} holds because $\Lambda$ is a multiple of the
identity and all leverages are equal:
$g_c\tp\Lambda g_c=\tfrac{2-\sqrt2}{2}\cdot2=2-\sqrt2=\cv$, which is
already \eqref{eq:quadric}.  For \eqref{eq:statg} take $c=0$.  The pairs
containing $0$ are $C_0$ and $C_1$, and
$\langle r_0,g_0\rangle=\langle r_1,g_0\rangle=\sqrt2$, so, since
normalizing $r_i$ twice contributes $\lVert r_i\rVert^{-2}$,
\[
  \sum_{i\,:\,0\in C_i}\lambda_i\langle u_i,g_0\rangle u_i
  =\tfrac14\cdot\frac{\sqrt2}{4+2\sqrt2}\,\bigl(r_0+r_1\bigr)
  =\tfrac14\cdot\frac{\sqrt2}{4+2\sqrt2}\,(2,0)
  =\frac{\sqrt2-1}{4}\,(1,0),
\]
because
$\dfrac{\sqrt2}{4+2\sqrt2}=\dfrac{\sqrt2\,(4-2\sqrt2)}{8}
=\dfrac{4\sqrt2-4}{8}=\dfrac{\sqrt2-1}{2}$, while
\[
  t_0\Lambda g_0=\tfrac14\cdot\tfrac{2-\sqrt2}{2}\,(\sqrt2,0)
  =\frac{2\sqrt2-2}{8}\,(1,0)=\frac{\sqrt2-1}{4}\,(1,0).
\]
The three remaining vectors are handled in the same way.  Finally, the
two pairs omitted from the active family are orthogonal:
$S_{\{0,3\}}=S_{\{1,2\}}=2I_2$, so $\lmin=2$ there, $F=2$, and
$S_{\{0,3\}}-I=I\pd0$.
\end{proof}

The datum of Theorem~\ref{thm:reach} is stationary for the maximum taken
over the four chosen pairs, which minimize $\lmin(\SC)$, and not for $F$.  It violates
\eqref{eq:argmax}, a disjunctive semialgebraic condition rather than an
equation and therefore beyond the algebra of \S\ref{sec:laws}.  It also
violates the hypothesis of every exclusion: its multiplier is isotropic
rather than of rank one, its four tight directions lie in $\R^2$ and
are dependent, and its active family is four distinct pairs, neither a
singleton nor a partition.

\subsection{The residue}\label{sec:census}

The residue is bounded below by \eqref{eq:count} and above by
Carath\'eodory's theorem applied to the convex combination
\eqref{eq:zeta}.

\begin{proposition}\label{prop:carath}
A stationarity datum produced by Theorem~\ref{thm:firstorder} may be
chosen with
\begin{equation}\label{eq:carath}
  \lvert A\rvert\;\le\;m(k+1)-k^{2} .
\end{equation}
At $(6,3)$ this is $15$, and at $(7,3)$ it is $19$.
\end{proposition}

\begin{proof}
After the spectral decomposition carried out in the proof of
Theorem~\ref{thm:firstorder}, identity \eqref{eq:zetarankone} reads
$\psi=\sum_{i\in A}\lambda_i\,\phi_{C_i,u_iu_i\tp}$, a convex
combination of the rank-one functionals of \eqref{eq:pieceSub}; and
$\psi$ vanishes on $\tang$.  So the origin lies in the convex hull of the
restrictions to $\tang$ of those functionals.  By \eqref{eq:rotationkill}
each restriction annihilates the subspace $\mathcal{O}$ of infinitesimal
rotations, which by \eqref{eq:orbit} lies in $\tang$ and has dimension
$\tfrac12k(k-1)$.  The combination therefore takes place in the
annihilator of $\mathcal{O}$ inside the dual of $\tang$, a space of dimension
\[
  \dim \tang-\dim\mathcal{O}
  =\Bigl(mk+m-\tfrac12k(k+1)-1\Bigr)-\tfrac12k(k-1)
  =mk+m-k^{2}-1 ,
\]
by Lemma~\ref{lem:tangent}.  Carath\'eodory's theorem
\cite{Caratheodory1911}, \cite[\S17]{Rockafellar1970} writes a point of a
convex hull in a space of dimension $d$ as a combination of at most $d+1$
of the points; here $d+1=mk+m-k^{2}=m(k+1)-k^2$.  Keep those indices,
discard the rest, and read the last two steps of the proof of
Theorem~\ref{thm:firstorder} again from the reduced combination: the
value, the subsets and the tight directions are unchanged, the
multipliers are the new ones, and $\Lambda$ is recomputed by
\eqref{eq:multiplier}.
\end{proof}

The dimension in that proof is the dimension of the design manifold
modulo the orthogonal group: $mk$ coordinates for the vectors and $m$ for
the weights, less the $\tfrac12k(k+1)+1$ constraints, less the
$\tfrac12k(k-1)$ rotations.  At $(6,3)$ it is $18+6-6-1-3=14$, and at
$(7,3)$ it is $21+7-6-1-3=18$.

Combining, at $(6,3)$ the active family covers $\{1,\dots,6\}$ and has
between $2$ and $15$ members, the value $2$ occurring only in the
partition branch of Corollary~\ref{cor:e3}; at $(7,3)$ it covers
$\{1,\dots,7\}$ and has between $3$ and $19$ members.  The number of
possibilities is finite, and at $(6,3)$ small enough to be listed.

\begin{proposition}\label{prop:census}
Up to relabelling, the families of $3$-subsets of $\{1,\dots,6\}$ whose
union is $\{1,\dots,6\}$ number $2102$, distributed by cardinality as in
Table~\ref{tab:census}.  Exactly one of them has two members, namely a
partition into two triples; so $2101$ have at least three members, and
$2069$ have at most fifteen.  Up to relabelling, the families of
$3$-subsets of $\{1,\dots,7\}$ whose union is $\{1,\dots,7\}$ number
$7\,011\,184$; every one has at least three members, and $5\,249\,381$
have at most nineteen.
\end{proposition}

The exclusions reach only a corner of this census.  Of the $2136$ classes of families of $3$-subsets
of $\{1,\dots,6\}$, Lemma~\ref{lem:overlap}(iv) removes the $34$ that do
not cover; Proposition~\ref{prop:saturated} removes $23$ of the
remainder, namely those in which some index lies in every member; and
Corollary~\ref{cor:e3} leaves the single partition class undecided.  So
$57$ classes are excluded and $2079$ are not.  The exclusions act at the
two ends of the range and are silent in the middle, where the census is
concentrated: they remove classes of cardinality at most three and
classes large enough to be saturated, while the distribution of
Table~\ref{tab:census} peaks at cardinality ten.  Conjecture~\ref{conj:sixthree}
is therefore not a finite verification away from the exclusions proved
here.

\begin{proof}
Orbit enumeration under the symmetric group.  For $n=6$ the $2^{20}$
families of triples were partitioned into orbits directly.  Independently,
for $n=6$ and $n=7$ the number of orbits of covering families of
cardinality $j$ was computed from Burnside's lemma as
\[
  \frac1{n!}\sum_{\varsigma\in S_n}\ \sum_{R}(-1)^{r-\lvert R\rvert}
    \bigl[x^{j}\bigr]\!\!\prod_{\kappa\,\subseteq\, T_R}
      \bigl(1+x^{\lvert\kappa\rvert}\bigr),
\]
where $\varsigma$ has $r$ cycles on $\{1,\dots,n\}$, the index $R$ runs over
sets of those cycles, $T_R$ is their union, and $\kappa$ runs over the
cycles of the permutation induced by $\varsigma$ on the $3$-subsets contained
in $T_R$.  A family fixed by $\varsigma$ is a union of such cycles, and its
union of points is $\varsigma$-invariant, so the alternating sum is
inclusion--exclusion over the sets a fixed family can avoid.  The two
computations agree at $n=6$.
\end{proof}

\begin{table}[ht]
\small
\centering
\begin{tabular}{@{}l*{10}{r}@{}}
\toprule
$j$   & 2 & 3 & 4 & 5 & 6 & 7 & 8 & 9 & 10 & 11\\
$N_j$ & 1 & 3 & 15 & 37 & 88 & 157 & 247 & 311 & 351 & 312\\
\midrule
$j$   & 12 & 13 & 14 & 15 & 16 & 17 & 18 & 19 & 20 & \\
$N_j$ & 249 & 161 & 94 & 43 & 21 & 7 & 3 & 1 & 1 & \\
\bottomrule
\end{tabular}
\medskip
\caption{The number $N_j$ of families of $j$ triples covering
$\{1,\dots,6\}$, up to relabelling.  The total is $2102$; the
Carath\'eodory bound \eqref{eq:carath} discards the last five columns,
leaving $2069$.}
\label{tab:census}
\end{table}

\begin{figure}[ht]
  \centering
  % The (6,3) census of Proposition prop:census, by cardinality of the
% active family.  Hatched: the window 3 <= |A^+| <= 15 left by
% Corollary cor:e3 (below) and Proposition prop:carath (above), which
% holds 2068 of the 2102 classes.
\begin{tikzpicture}
\begin{axis}[
  width=13.4cm, height=6.6cm,
  scale only axis,
  ybar, bar width=4.3mm,
  xmin=1.3, xmax=20.7,
  ymin=0, ymax=470,
  xtick={2,3,4,5,6,7,8,9,10,11,12,13,14,15,16,17,18,19,20},
  ytick={0,100,200,300},
  xlabel={cardinality $\lvert A^{+}\rvert$ of the active family},
  ylabel={classes under $S_6$},
  xlabel style={font=\small}, ylabel style={font=\small},
  tick label style={font=\scriptsize},
  axis lines=left,
  every axis plot/.append style={bar shift=0pt},
  nodes near coords,
  nodes near coords style={font=\tiny, rotate=90, anchor=west, inner sep=1.5pt},
  clip=false]

% ---- outside the window, below --------------------------------------
\addplot[draw=black, line width=0.4pt, fill=black!12]
  coordinates {(2,1)};

% ---- the window 3 <= |A| <= 15 --------------------------------------
\addplot[draw=black, line width=0.4pt,
         pattern=north east lines, pattern color=black!60]
  coordinates {(3,3) (4,15) (5,37) (6,88) (7,157) (8,247) (9,311)
               (10,351) (11,312) (12,249) (13,161) (14,94) (15,43)};

% ---- outside the window, above --------------------------------------
\addplot[draw=black, line width=0.4pt, fill=black!12]
  coordinates {(16,21) (17,7) (18,3) (19,1) (20,1)};

% ---- the cut-offs ----------------------------------------------------
\draw[dashed, line width=0.7pt] (axis cs:2.5,0) -- (axis cs:2.5,395);
\draw[dashed, line width=0.7pt] (axis cs:15.5,0) -- (axis cs:15.5,395);
\node[font=\scriptsize, anchor=north west, align=left]
  at (axis cs:1.7,462)
  {below: coverage and Corollary~\ref{cor:e3} force
   $\lvert A^{+}\rvert\ge3$,\\
   except in the single partition class at $\lvert A^{+}\rvert=2$};
\node[font=\scriptsize, anchor=north east, align=right]
  at (axis cs:20.4,462)
  {above: the Carath\'eodory bound \eqref{eq:carath} is $15$\\
   $2102$ classes in all; the window holds $2068$};
\end{axis}
\end{tikzpicture}

  \caption{The census of Proposition~\ref{prop:census} at $(6,3)$.  The
  active family covers $\{1,\dots,6\}$, so it has at least two members,
  and the one class of cardinality two is the partition that
  Corollary~\ref{cor:e3} leaves undecided; \eqref{eq:carath} removes the
  classes above fifteen.}
  \label{fig:census}
\end{figure}

\begin{conjecture}\label{conj:sixthree}
No design at $(6,3)$ carries an admissible stationarity datum with
$0<\cv<1$.
\end{conjecture}

\begin{conjecture}\label{conj:seventhree}
No design at $(7,3)$ carries an admissible stationarity datum with
$0<\cv<1$.
\end{conjecture}

Admissibility may not be dropped: Theorem~\ref{thm:reach} refutes the
corresponding statement at $(4,2)$ without it.  With it, each conjecture
is a finite conjunction of semialgebraic emptiness questions, one for
each of the classes counted in Proposition~\ref{prop:census}.  For a
fixed active family the unknowns are the $mk$ coordinates, the $m$
weights, the multipliers, the tight directions, the entries of $\Lambda$
and the value $\cv$; the conditions are \eqref{eq:statg},
\eqref{eq:statt}, $\lVert u_i\rVert=1$, $S_{C_i}u_i=\cv u_i$,
$\lambda_i\ge0$, $\sum_i\lambda_i=1$, \eqref{eq:parseval}, $t\pd0$,
$\sum_ct_c=1$, $0<\cv<1$, and \eqref{eq:argmax} --- the last a disjunction
of sign conditions on the leading principal minors of $\SC-\cv I$, one
disjunction for each $k$-subset $C$.

\begin{theorem}[conditional]\label{thm:rankthree}
Assume Conjectures~\ref{conj:sixthree} and~\ref{conj:seventhree}, and
assume that a failure of $\Gtz(6,3)$, respectively of $\Gtz(7,3)$,
produces an extremal counterexample at that cell in the sense of
Definition~\ref{def:extremalcex}.  Then $\Gtz(m,3)$ holds for every $m$.
\end{theorem}

\begin{proof}
By Corollary~\ref{cor:twocells} it suffices to prove $\Gtz(6,3)$ and
$\Gtz(7,3)$.  Suppose $\Gtz(6,3)$ fails.  By hypothesis $F$ attains a
local minimum over the designs at $(6,3)$ at a design of value $\cv<1$,
and $\cv>0$ by Lemma~\ref{lem:fpos}.  Theorem~\ref{thm:firstorder}
produces at that design a stationarity datum with value $\cv$, and it is
admissible because $\cv=F$ there.  This contradicts
Conjecture~\ref{conj:sixthree}.  The same argument at $(7,3)$ with
Conjecture~\ref{conj:seventhree} gives $\Gtz(7,3)$.
\end{proof}

The extremality hypothesis is all this section takes from
\S\ref{sec:compactness}; Theorem~\ref{thm:variational} supplies it for
the chart objective.

At $(6,3)$ the conjecture is an enumeration; at $(7,3)$ it is not, and an
argument uniform in the active family is required.  The exclusions of \S\ref{sec:exclusions} remove the small and the
transversal configurations, and Example~\ref{ex:splitseven} shows where the residue
lives: at the frontier size the system is realized, at the threshold, by
an active family of twenty triples.

\begin{wip}
Decide the partition branch of Corollary~\ref{cor:e3}: either exclude a
datum whose active family is a two-block partition of $\{1,\dots,6\}$
with $0<\cv<1$, or exhibit one.
\end{wip}
